\section{5. Reações químicas}

\begin{frame}{Obtenção do silício a partir da areia}
  \begin{block}{Forno a arco elétrico, \grau{\sim 2000}}
    \[ \qui{SiO_2}\,(s) + 2\,\qui{C}\,(s) \seta \qui{Si}\,(l) + 2\,\qui{CO}\,(g) \]
  \end{block}
  \vspace{0.2cm}
  \begin{itemize}
    \item O carbono tira o oxigênio da areia, numa reação de
          \termo{oxirredução}.
    \item Gasta muita energia, e por isso a produção fica onde a energia é
          barata.
    \item O silício sai com 98\% de pureza, ainda \alert{muito longe} do que
          a eletrônica precisa.
  \end{itemize}
\end{frame}

\begin{frame}{Purificação pelo processo Siemens}
  \begin{block}{1. O silício vira gás, a \grau{300}}
    \[ \qui{Si}\,(s) + 3\,\qui{HCl}\,(g) \seta \qui{SiHCl_3}\,(g) + \qui{H_2}\,(g) \]
  \end{block}
  \begin{block}{2. O gás volta a ser silício, a \grau{1100}}
    \[ \qui{SiHCl_3}\,(g) + \qui{H_2}\,(g) \seta \qui{Si}\,(s) + 3\,\qui{HCl}\,(g) \]
  \end{block}
  \vspace{0.2cm}
  \centering\small
  Purificar sólido é difícil, mas o gás pode ser \termo{destilado} como se
  destila álcool. A segunda reação é a primeira \alert{ao contrário}, e o
  silício volta sem a sujeira.
\end{frame}

\begin{frame}{Reações usadas na fabricação}
  \begin{block}{Cria a camada isolante sobre o silício}
    \[ \qui{Si}\,(s) + \qui{O_2}\,(g) \seta \qui{SiO_2}\,(s) \]
  \end{block}
  \vspace{0.2cm}
  \begin{block}{Forma o cristal de arseneto de gálio}
    \[ \qui{Ga(CH_3)_3}\,(g) + \qui{AsH_3}\,(g) \seta \qui{GaAs}\,(s) + 3\,\qui{CH_4}\,(g) \]
  \end{block}
  \vspace{0.3cm}
  \centering\small
  Para desenhar o circuito, o ácido fluorídrico abre janelas na camada
  isolante sem atacar o silício.
\end{frame}
